\documentclass[sigconf,review,anonymous]{acmart}

\makeatletter
\global\@ACM@balancefalse
\makeatother
\raggedbottom
\settopmatter{printacmref=true}
\renewcommand\footnotetextcopyrightpermission[1]{}

\acmConference[CIKM '26]{35th ACM International Conference on Information and Knowledge Management}{November 2026}{Rome, Italy}
\acmYear{2026}
\copyrightyear{2026}
\acmDOI{}
\acmISBN{}

\title{Lag-Aware Self-Supervised Dynamic Graph Structure Discovery for Multivariate Time Series}

\author{Anonymous Authors}
\affiliation{%
  \institution{Anonymous Institution}
  \city{Anonymous City}
  \country{Anonymous Country}}
\email{anonymous@example.com}

\begin{document}

\begin{abstract}
Discovering directed and lag-aware structure from multivariate temporal signals is central to interpretable forecasting, anomaly analysis, and dynamic graph modeling, yet real-world edge supervision is rarely available.
This paper studies whether self-supervised temporal reconstruction signals can recover sparse directed lagged graph structure without edge labels.
We develop STCG, a lag-aware graph-state discovery framework that scores local temporal dependencies in sliding windows, clusters windows into latent graph states, and aggregates state-specific edge scores.
On synthetic systems with known directed lagged ground truth, the graph-state variant improves switching-regime recovery over dynamic VAR, dynamic Granger, and PCMCI-ParCorr baselines while exposing interpretable latent states aligned with hidden regimes.
The current evidence supports a focused claim: graph-state consistency helps recover changing temporal dependency structure, while stationary and short-sequence regimes remain important limitations.
\end{abstract}

\begin{CCSXML}
<ccs2012>
   <concept>
       <concept_id>10002951.10003227.10003236</concept_id>
       <concept_desc>Information systems~Data mining</concept_desc>
       <concept_significance>500</concept_significance>
   </concept>
   <concept>
       <concept_id>10010147.10010257.10010258.10010261</concept_id>
       <concept_desc>Computing methodologies~Machine learning approaches</concept_desc>
       <concept_significance>500</concept_significance>
   </concept>
</ccs2012>
\end{CCSXML}

\ccsdesc[500]{Information systems~Data mining}
\ccsdesc[500]{Computing methodologies~Machine learning approaches}

\keywords{dynamic graph learning, self-supervised learning, temporal causal discovery, multivariate time series}

\maketitle

\section{Introduction}

Many multivariate time series are generated by sparse lagged interactions, but the dependency graph is rarely observed.
Forecasting models can exploit such dependencies implicitly, yet the learned representation is often difficult to audit as a directed, lag-aware graph.
This creates a gap between predictive temporal modeling and interpretable structure discovery.

We focus on the unsupervised setting: only node time series are observed, and no edge labels are provided during learning.
The central question is whether reconstruction-derived temporal signals can recover useful sparse directed lagged edges, especially when the active graph changes across regimes.
Our current answer is deliberately scoped.
On switching synthetic systems, a graph-state-aware self-supervised scorer improves lagged edge recovery over independent sliding-window baselines.
On stationary systems, global VAR, Granger, and PCMCI baselines remain stronger, suggesting that dynamic graph-state modeling should be used when graph changes are plausible rather than as a universal replacement.
This paper makes three focused contributions.
First, it formulates lag-aware graph discovery as label-free sparse temporal reconstruction over sliding windows.
Second, it introduces latent graph-state selection and state-level aggregation as a way to denoise changing local dependency estimates.
Third, it provides a bounded evaluation protocol that separates synthetic structure recovery from real-data forecasting, state-diagnosis, and transition-warning probes.

\section{Related Work}

Classical temporal discovery begins with Granger-style predictive tests, which ask whether lagged values of one variable improve prediction of another~\cite{granger1969investigating}.
Modern conditional-independence approaches, including PCMCI-style temporal discovery, broaden this view to larger systems with autocorrelation and indirect links~\cite{runge2019detecting}.
These methods are strong references for causal discovery; we therefore include a PCMCI-ParCorr baseline in the synthetic recovery loop.
They are usually evaluated as fixed conditional-dependence procedures over a series or a set of regimes rather than as latent dynamic graph-state learners.
In parallel, graph-based temporal forecasting models learn or exploit graph structure for prediction~\cite{yu2018spatio,wu2020connecting}, and Transformer-style time-series models provide powerful sequence forecasting baselines~\cite{zhou2021informer}.
Their objective is typically forecasting accuracy, and learned graph or attention weights are often auxiliary model internals.
Our primary question is narrower: whether reconstruction-derived signals recover directed lagged structure when the graph changes across regimes.
This distinction motivates a structure-recovery-first evaluation: real forecasting and warning tasks are treated as diagnostic probes, not as claims of forecasting state of the art.

\section{Problem Formulation}

Given multivariate observations $X_{1:T} \in \mathbb{R}^{T \times N}$, the goal is to learn latent graphs whose edges can be directed and associated with a discrete lag.
No observed edge labels are used during training.
Evaluation on synthetic data compares recovered edges and lags against known ground truth.
For a maximum lag $L$, the recovered score tensor $S \in \mathbb{R}_{\geq 0}^{L \times N \times N}$ uses $S_{\ell i j}$ to score the directed relation $i_{t-\ell}\rightarrow j_t$.
Self-edges are excluded.

\section{Method}

\subsection{Lag-Aware Window Scoring}

For each ordered pair $(i,j)$ and lag $\ell \in \{1,\ldots,L\}$, the method estimates a nonnegative score $s_{ij\ell,w}$ inside window $w$.
The current implementation uses sparse lagged reconstruction as the base self-supervised signal: a lagged design matrix predicts each target node from past source-node values, and coefficient magnitudes become candidate edge scores.
Self-edges are masked.
For window $w=[a,b)$ and target node $j$, define a lagged design matrix $Z_w$ with columns $X_{t-\ell,i}$ and target vector $y_{w,j}=X_{t,j}$.
STCG-v1 fits
\[
\hat{\beta}_{w,j}
= \arg\min_{\beta}
\|y_{w,j}-Z_w\beta\|_2^2 + \lambda \|\beta\|_1,
\]
and maps coefficients to scores by $S_{w,\ell i j}=|\hat{\beta}_{w,j,\ell i}|$.
The coefficient estimator is intentionally close to a dynamic VAR-Lasso baseline.
The distinguishing step is what happens after local scoring: instead of treating each window independently or averaging all windows, STCG-v1 treats window-level score tensors as observations of latent graph regimes.

\subsection{Latent Graph-State Aggregation}

Independent window estimates are noisy, and max aggregation alone can overfit window-specific fluctuations.
STCG-v1 therefore clusters windows using flattened non-self lagged edge-score features.
For each inferred graph state, it aggregates member-window scores, then combines state-level graphs by a max operator to recover edges that are active in at least one regime.
The clustering uses only reconstruction-derived score features, not synthetic edge labels or regime labels.
Let $z_w=\mathrm{vec}_{i\neq j}(S_w)$ be the non-self edge-score feature for window $w$.
K-means assigns a latent state $c_w \in \{1,\ldots,K\}$.
The automatic state gate accepts a multi-state partition when its K-means inertia improvement exceeds a threshold, or when a lower-improvement partition is temporally coherent and contains no tiny state.
For the latter case, the implementation uses an adjusted temporal coherence score that compares adjacent-window state persistence with the persistence expected from the inferred state frequencies.
Each state graph is
\[
G_k = \mathrm{Agg}\{S_w: c_w=k\},
\]
where the default aggregator is the element-wise mean.
The final score tensor is
\[
S_{\ell i j} = \max_k G_{k,\ell i j}.
\]
The label-free state selector scores each candidate state count $K$ by
\[
q_K = \frac{\Delta_K \tau_K b_K}{\log_2(K+1)},
\]
where $\Delta_K$ is the score-clustering improvement over one state, $\tau_K$ is adjusted temporal coherence, and $b_K=\min(1,K\pi_{\min})$ penalizes tiny states with minimum state fraction $\pi_{\min}$.
For the Hydraulic transition-warning ablation, we also test a transition-oriented score
\[
q^{\mathrm{warn}}_K =
\frac{\Delta_K b_K(1+\rho_K)/2}{\log_2(K+1)}, \quad
\rho_K = 1-\min(1,\tau_K),
\]
which still requires a balanced multi-state partition but does not reject low-coherence state changes that may be useful for warning.

\subsection{Algorithm}

\begin{figure}[t]
\centering
\small
\begin{tabular}{p{0.96\linewidth}}
\toprule
\textbf{Algorithm 1: STCG-v1 graph-state scoring} \\
\textbf{Input:} time series $X_{1:T}$, maximum lag $L$, window size $h$, stride $r$, candidate state counts $\mathcal{K}$. \\
1. Form overlapping windows $w_1,\ldots,w_M$ using size $h$ and stride $r$. \\
2. For each window and target node, fit sparse lagged reconstruction and store non-self score tensor $S_w$. \\
3. Flatten each $S_w$ into a non-self graph feature $z_w$. \\
4. For each $K\in\mathcal{K}$, cluster $\{z_w\}$ into graph states and compute the label-free score $q_K$. \\
5. Select the highest-scoring valid $K$; if no multi-state candidate passes the gate, use one state. \\
6. Aggregate window score tensors within each selected state to obtain $G_k$. \\
7. Return $S_{\ell i j}=\max_k G_{k,\ell i j}$ for structure recovery, graph-prior forecasting, or warning probes. \\
\bottomrule
\end{tabular}
\caption{Implemented STCG-v1 pipeline. Synthetic edge labels, hidden regimes, and real external labels are never used in steps 1--7.}
\Description{Pseudocode listing the seven implemented STCG-v1 steps from sliding-window lagged reconstruction through graph-state clustering and state-level score aggregation.}
\label{alg:stcg-v1}
\end{figure}

\subsection{Graph-Guided Forecasting Probe}

For real-data probing, graph scores are converted either to a hard non-self feature mask or to a soft ridge prior.
The hard probe keeps the top-scoring non-self lagged sources for each target.
The soft probe keeps all lagged features and changes only the ridge penalty.
Given a target-specific feature weight $w_{mj}>0$, it fits
\[
\hat{\beta}_j
= \arg\min_{\beta}
\|y_j-Z\beta\|_2^2 + \alpha \sum_m \frac{\beta_m^2}{w_{mj}}.
\]
Weights are normalized from the score tensor per target, with self-lag features assigned unit weight.
Thus a high graph score weakens shrinkage for that lagged source, while low-score features remain available rather than being deleted.

\subsection{Real-Data Diagnostic Protocols}

Real benchmarks use chronological splits.
Graph scores, scalers, masks, and ridge priors are fit only on the training segment before one-step forecasting is evaluated on future samples.
For state-label diagnostics, external labels are used after unsupervised graph-state inference only to map inferred states or graph centroids to labels on a calibration subset; held-out windows report balanced accuracy and macro F1.
For Hydraulic transition warnings, each example uses only a lookback of past cycles and the target is whether the audited label changes within the next horizon.
Classifiers are trained on the earliest segment, thresholds and optional STCG warning configurations are selected on the following chronological validation segment, and final AP/F1/ROC AUC are reported once on the final test segment.

\subsection{Prototype Variants}

STCG-v0 trains differentiable lagged gates from temporal reconstruction, but it does not yet outperform the best dynamic VAR baseline.
STCG-v1 keeps the self-supervised reconstruction signal but adds latent graph-state consistency.
An optional state-contrast term was tested and removed from the default because it slightly reduced switching-regime recovery.
Thus STCG-v1 should not be read as only DynVAR followed by a post-hoc clustering plot.
Dynamic VAR baselines estimate local graphs independently and report those window scores directly.
STCG-v1 instead uses local reconstruction scores as observations of unobserved graph regimes, selects the state count without labels, aggregates graphs within inferred states, and returns both a recovered lagged graph and an interpretable state sequence.
This is the mechanism tested by the no-state and state-only ablations.

\section{Experiments}

\subsection{Synthetic Structure Recovery}

The primary experiments use VAR, nonlinear VAR, and switching-regime generators with known directed edges and lags.
All reported synthetic tables use twenty random seeds.
Metrics are Edge AUC over lagged directed edges, Pair AUC after collapsing lags, Precision@K at the number of true edges, Structural Hamming Distance (SHD), and lag recovery accuracy for correctly recovered directed pairs.
Paired tests compare methods on the same seeds, with bootstrap confidence intervals, one-sided paired $t$-tests, and one-sided Wilcoxon signed-rank tests used as supporting evidence rather than a substitute for effect-size inspection.

\begin{table}[t]
\caption{Switching-regime structure recovery over twenty seeds. Higher is better for Edge AUC, Pair AUC, Precision@K, and lag accuracy; lower is better for SHD.}
\label{tab:switching-recovery}
\centering
\small
\begin{tabular}{lrrrrr}
\toprule
Method & Edge AUC & Pair AUC & P@K & SHD & Lag Acc. \\
\midrule
VAR-Lasso & 0.780 & 0.808 & 0.625 & 40.8 & 0.767 \\
Granger-F & 0.778 & 0.806 & 0.622 & 41.2 & 0.770 \\
PCMCI-ParCorr & 0.779 & 0.809 & 0.621 & 41.3 & 0.760 \\
DynVAR-Lasso & 0.798 & 0.830 & 0.646 & 38.6 & 0.797 \\
DynGranger-F & 0.797 & 0.830 & 0.650 & 38.1 & 0.791 \\
STCG-v0 & 0.780 & 0.790 & 0.615 & 41.9 & 0.777 \\
STCG-v1 & \textbf{0.812} & \textbf{0.839} & \textbf{0.673} & \textbf{35.5} & \textbf{0.811} \\
\bottomrule
\end{tabular}
\Description{Synthetic switching-regime recovery table comparing VAR, Granger, dynamic VAR, dynamic Granger, STCG-v0, and STCG-v1 by edge AUC, pair AUC, precision at K, structural Hamming distance, and lag accuracy.}
\end{table}

Table~\ref{tab:switching-recovery} shows that the graph-state variant improves over dynamic VAR, dynamic Granger, and PCMCI-ParCorr baselines on switching-regime recovery.
Across all generator types, STCG-v1 is strongest on switching VAR, while global VAR-Lasso, Granger-F, and PCMCI-ParCorr are stronger on stationary and nonlinear VAR.
This supports a dynamic-structure claim rather than a blanket claim over all temporal systems.
Paired tests over the same twenty switching seeds show significant gains for Edge AUC, Precision@K, and SHD under both paired $t$ and Wilcoxon tests; Table~\ref{tab:paired-switching} gives the effect sizes, bootstrap intervals, and both p-values.
Pair AUC and lag accuracy improve in mean but are not statistically significant at this sample size.
Against DynGranger-F, STCG-v1 likewise improves switching Edge AUC, Precision@K, and SHD under both tests; lag accuracy improves under the paired $t$-test but is weaker under Wilcoxon, and Pair AUC again improves only in mean.
Table~\ref{tab:paired-switching} reports the full DynVAR-Lasso paired comparison, and the DynGranger-F paired table is included in the reproduction artifacts.
We also add a focused PCMCI-ParCorr comparison using Tigramite with ParCorr conditional tests.
This is a strong stationary causal-discovery reference: on stationary VAR, PCMCI-ParCorr reaches Edge AUC $0.920$ versus STCG-v1's $0.876$, and on nonlinear VAR it reaches $0.984$ versus $0.969$.
On switching VAR, however, STCG-v1 improves over PCMCI-ParCorr on all five metrics; the paired Edge AUC gain is $0.033$ with one-sided paired $t$-test $p=0.0005$ and Wilcoxon $p=0.0004$.
To probe the mechanism, we add an oracle-state diagnostic that uses true synthetic regimes only after local window scores have been computed.
On the same twenty switching seeds, STCG-v1 inferred-state aggregation nearly matches this oracle upper bound: Edge AUC is $0.812$ versus oracle $0.813$ and DynVAR-Lasso/no-state $0.798$, while Precision@K is $0.673$ versus oracle $0.672$ and DynVAR-Lasso/no-state $0.646$.
Thus, under the default switching setup, most of the available oracle-state gain is already captured by the label-free graph-state inference.

\begin{table}[t]
\caption{Paired STCG-v1 versus DynVAR-Lasso tests on switching-regime recovery. Positive improvement means STCG-v1 is better; for SHD it means fewer errors. $p_t$ is the one-sided paired $t$-test and $p_W$ is the one-sided Wilcoxon signed-rank test.}
\label{tab:paired-switching}
\centering
\small
\resizebox{\linewidth}{!}{\begin{tabular}{lrrrlrrl}
\toprule
Metric & STCG-v1 & DynVAR & Improve & 95\% CI & $p_t$ & $p_W$ & W/T/L \\
\midrule
Edge AUC & 0.812 & 0.798 & 0.0143 & [0.0067, 0.0213] & 0.0007 & 0.0014 & 17/0/3 \\
Pair AUC & 0.839 & 0.830 & 0.0092 & [-0.0081, 0.0272] & 0.1683 & 0.2045 & 11/0/9 \\
P@K & 0.673 & 0.646 & 0.0274 & [0.0150, 0.0403] & 0.0003 & 0.0012 & 14/4/2 \\
SHD & 35.5 & 38.6 & 3.1 & [1.7, 4.6] & 0.0003 & 0.0007 & 14/4/2 \\
Lag Acc & 0.811 & 0.797 & 0.0143 & [-0.0055, 0.0366] & 0.1054 & 0.1896 & 9/6/5 \\
\bottomrule
\end{tabular}
}
\Description{Paired synthetic comparison table showing STCG-v1, DynVAR-Lasso, mean improvement, bootstrap confidence intervals, paired t-test p-values, Wilcoxon p-values, and win-tie-loss counts.}
\end{table}

\begin{table}[t]
\caption{Paired STCG-v1 versus PCMCI-ParCorr tests on switching-regime recovery. Positive improvement means STCG-v1 is better; for SHD it means fewer errors.}
\label{tab:paired-pcmci}
\centering
\small
\resizebox{\linewidth}{!}{\begin{tabular}{lrrrlrrl}
\toprule
Metric & STCG-v1 & PCMCI-ParCorr & Improve & 95\% CI & $p_t$ & $p_W$ & W/T/L \\
\midrule
Edge AUC & 0.812 & 0.779 & 0.0334 & [0.0161, 0.0492] & 0.0005 & 0.0004 & 17/0/3 \\
Pair AUC & 0.839 & 0.809 & 0.0294 & [0.0099, 0.0484] & 0.0043 & 0.0016 & 17/0/3 \\
P@K & 0.673 & 0.621 & 0.0520 & [0.0296, 0.0709] & $<10^{-4}$ & 0.0007 & 18/1/1 \\
SHD & 35.5 & 41.3 & 5.8 & [3.3, 7.9] & $<10^{-4}$ & 0.0006 & 18/1/1 \\
Lag Acc & 0.811 & 0.760 & 0.0511 & [0.0240, 0.0775] & 0.0008 & 0.0028 & 15/3/2 \\
\bottomrule
\end{tabular}
}
\Description{Paired synthetic comparison table showing STCG-v1, PCMCI-ParCorr, mean improvement, bootstrap confidence intervals, paired t-test p-values, Wilcoxon p-values, and win-tie-loss counts.}
\end{table}

\subsection{Real Benchmarks}

The repository now includes a leakage-free CSV forecasting harness: graph scores are fit on the training segment only, then one-step forecasting, top learned edges, and inferred state windows are exported.
We separate real-data evidence into three probes: one-step forecasting, graph-state diagnosis against external operating labels, and future transition warning.
As initial real benchmarks, we evaluate four ETT files~\cite{zhou2021informer}, BeijingPM25 from the Beijing multi-site air-quality data~\cite{zhang2017cautionary}, OccupancySensors from UCI Occupancy Detection~\cite{candanedo2016occupancydata,candanedo2016accurate}, AReMActivities from UCI AReM~\cite{palumbo2016arem}, UCIHARWindowSignals from UCI smartphone activity recognition~\cite{reyesortiz2012har}, and HydraulicSystems from UCI hydraulic condition monitoring~\cite{helwig2018hydraulic}.
They cover 17,420--69,680 ETT samples, 35,064 hourly PM2.5 observations, 20,560 minute-level room-sensor observations, 42,237 ordered AReM wearable-sensor samples, 10,299 UCI HAR inertial-signal windows, and 132,300 one-second Hydraulic Systems sensor rows from 2,205 test cycles.
Table~\ref{tab:real-forecasting} reports one-step forecasting RMSE after a chronological train/test split.
Dense ridge is strongest on the ETT, BeijingPM25, and AReMActivities default runs among materially distinct models; on OccupancySensors, persistence is slightly best by $1.2\times 10^{-4}$ RMSE, below the reporting threshold used for meaningful real-data gains.
STCG-v1 graph masking ties the DynVAR-Lasso graph mask and underperforms dense ridge and persistence.
The learned top edges are interpretable load/temperature-channel, station-to-station, room-sensor, activity-sensor, or hydraulic-sensor relations.
Default graph-state diagnostics still collapse on most non-mode-labelled real datasets and on UCI HAR, but the temporal fallback recovers non-collapsed states on AReMActivities and HydraulicSystems.
On OccupancySensors, post-hoc occupancy alignment gives ARI $0.000$ and aligned accuracy $0.737$, equal to the majority window-label rate.
On AReMActivities, the default $K=2$ activity alignment is non-collapsed but moderate (ARI $0.251$); a focused threshold sweep shows $K=3$ reaches ARI $0.309 \pm 0.035$, while $K=4$ gives the best weighted-forecasting RMSE.
An additional hard-mask sweep over graph fraction, Lasso strength, and window size does not find any STCG-v1 mask that beats dense ridge or persistence; the non-self dense baseline also underperforms, indicating that self-lag features are central for these forecasting splits.
Graph-weighted ridge avoids this deletion effect: a sweep yields a meaningful AReMActivities gain of $0.0029$ RMSE over tuned dense ridge, while all other datasets remain below the $0.001$ threshold.
Under tuned AReM $K=4$ states, STCG-v1 further improves weighted RMSE from DynVAR-Lasso's $0.9062$ to $0.9041$, but the five-seed paired test is not yet significant ($p=0.142$).
A second mode-labelled UCI HAR transfer check does not reproduce the AReM signal: non-collapsed settings reach only weak activity alignment (best ARI $0.090$) and do not meaningfully beat DynVAR-Lasso.
Hydraulic Systems gives a stronger industrial state-transfer result: default automatic $K=2$ graph states align with the cooler operating condition (ARI $0.735 \pm 0.006$, aligned accuracy $0.902 \pm 0.002$); this structural alignment still does not produce a meaningful weighted-RMSE gain over DynVAR-Lasso.
We therefore add a label-stratified state-diagnosis harness beyond one-step forecasting.
After unsupervised graph-state inference, a subset of window labels is used only to map states or centroids to external labels, and held-out windows are evaluated.
On Hydraulic cooler-condition labels, default $K=2$ compact state mapping underfits the three-level label (balanced accuracy $0.656 \pm 0.002$), while $K=3$ state mapping reaches $0.948 \pm 0.004$; graph-score centroid features reach $0.959 \pm 0.003$, and a simple sensor-feature centroid ceiling reaches $0.994 \pm 0.004$.
To avoid choosing $K$ from labels, we add a label-free selector that maximizes state improvement, temporal coherence, balance, and a logarithmic parsimony penalty.
It selects $K=3$ on Hydraulic in all five seeds, preserving $0.948 \pm 0.004$ balanced accuracy.
The same selector chooses coarse $K=2$ AReM states in all five seeds, reaching only $0.283 \pm 0.005$ balanced accuracy against seven activity labels, while graph-score centroids at $K=4$ reach $0.623 \pm 0.033$.
To reduce transductive-state concerns, we also add a forward-only Hydraulic check that fits graph-state centroids and state-label mappings on chronological train windows, then assigns future windows to the train-fitted centroids.
For cooler-condition late-future windows, STCG-v1 state mapping reaches balanced accuracy $0.953 \pm 0.003$; graph-score centroids reach $0.966$, and the sensor-feature ceiling reaches $1.000$.
This result is encouraging but narrow because the future cooler test segment contains only the late-life cooler condition; a multi-class valve-condition forward check is negative for compact STCG states (balanced accuracy $0.250$), while graph-score centroids reach $0.323$.
Finally, we add a chronological transition-warning pilot on Hydraulic labels with features computed only from past cycles and warning thresholds selected on a validation segment.
For 10-cycle valve-condition warnings, sensor-plus-graph logistic regression reaches AP $0.909$ and ROC AUC $0.771$ versus majority AP $0.742$; for 10-cycle stable-flag warnings, graph-only Ridge VAR features reach AP $0.269$ versus majority AP $0.244$.
Adding label-free selected STCG-v1 graph-state features with a transition-oriented selector and explicit state-dynamics features gives a mixed result: sensor-plus-STCG is best by AP on 5-cycle stable-flag warnings (AP $0.165$ versus majority AP $0.128$) and 3-cycle valve-condition warnings (AP $0.239$ versus majority AP $0.222$), and gives the best F1 on 10-cycle valve-condition warnings ($0.856$), but Ridge graph features remain stronger on the main 10-cycle valve-condition AP result.
An additional validation-selected warning ablation over stable versus transition selectors and small lookback/window choices improves sensor-plus-STCG on 10-cycle stable-flag AP to $0.302$, but the same lookback has a stronger sensor-only AP of $0.374$; on 10-cycle valve-condition, selected sensor-plus-STCG reaches AP $0.840$, still below sensor-plus-graph AP $0.909$.
Shorter warning horizons remain weak.

\begin{table}[t]
\caption{Initial real benchmark one-step forecasting RMSE on ETT small, BeijingPM25, OccupancySensors, AReMActivities, UCIHARWindowSignals, and HydraulicSystems datasets. Mask columns use hard graph feature selection; weight columns use graph-weighted ridge. Lower is better.}
\label{tab:real-forecasting}
\centering
\small
\resizebox{\linewidth}{!}{\begin{tabular}{lrrrrrr}
\toprule
Dataset & Persistence & Dense & DynVAR Mask & DynVAR Weight & STCG Mask & STCG Weight \\
\midrule
AReMActivities & 1.128 & 0.912 & 1.259 & 0.912 & 1.230 & 0.912 \\
BeijingPM25 & 0.266 & 0.233 & 0.367 & 0.233 & 0.367 & 0.233 \\
ETTh1 & 0.439 & 0.409 & 0.728 & 0.409 & 0.728 & 0.409 \\
ETTh2 & 0.246 & 0.227 & 0.888 & 0.227 & 0.888 & 0.227 \\
ETTm1 & 0.253 & 0.242 & 0.695 & 0.242 & 0.695 & 0.242 \\
ETTm2 & 0.156 & 0.149 & 0.802 & 0.149 & 0.802 & 0.149 \\
HydraulicSystems & 0.167 & 0.155 & 0.472 & 0.156 & 0.470 & 0.156 \\
OccupancySensors & 0.046 & 0.046 & 0.744 & 0.046 & 0.744 & 0.046 \\
UCIHARWindowSignals & 0.852 & 0.680 & 0.811 & 0.680 & 0.811 & 0.680 \\
\bottomrule
\end{tabular}
}
\Description{Real benchmark forecasting table reporting RMSE for persistence, dense ridge, DynVAR graph masks and weights, and STCG graph masks and weights across nine datasets.}
\end{table}

Table~\ref{tab:evidence-summary} summarizes how the main results should be read.
\begin{table}[t]
\caption{Evidence summary across synthetic and real probes.}
\label{tab:evidence-summary}
\centering
\small
\resizebox{\linewidth}{!}{\begin{tabular}{lll}
\toprule
Probe & Evidence & Reading \\
\midrule
Switching synthetic & AUC/P@K/SHD significant vs DynVAR, DynGranger, and PCMCI & Main positive claim \\
Stationary/nonlinear & VAR, Granger, or PCMCI stronger & Boundary condition \\
Real forecasting & Dense ridge or persistence usually strongest & Diagnostic only \\
State diagnosis & Hydraulic BA $0.948 \pm 0.004$; AReM moderate; HAR weak & Conditional transfer \\
Transition warning & STCG helps selected horizons, not main AP settings & Bounded pilot \\
\bottomrule
\end{tabular}
}
\Description{Summary table grouping the evidence into switching synthetic recovery, non-switching recovery, real forecasting, graph-state diagnosis, and transition warning, with the interpretation for each probe.}
\end{table}

\subsection{Baselines and Ablations}

Implemented synthetic baselines include random scores, lag correlation, global VAR-Ridge, VAR-Lasso, Granger-F, PCMCI-ParCorr, and dynamic window versions of the VAR and Granger families.
The real forecasting probe adds dense ridge, persistence, hard graph-mask ridge, and soft graph-weighted ridge.
The current benchmark does not yet include neural temporal graph forecasting baselines; this is treated as a limitation rather than hidden behind the main tables.
This omission limits broader forecasting claims, but it does not change the paired comparison against the implemented dynamic VAR, dynamic Granger, and PCMCI-ParCorr families.
Focused ablations test whether the switching-regime gain comes from ordinary windowing, latent graph-state aggregation, or the optional state-contrast term.

\begin{table}[t]
\caption{Focused ablation on twenty-seed switching-regime recovery. The gain comes from latent graph-state clustering rather than windowing alone.}
\label{tab:state-ablation}
\centering
\small
\begin{tabular}{lrrr}
\toprule
Variant & Edge AUC & P@K & SHD \\
\midrule
DynVAR-Lasso & 0.798 & 0.646 & 38.6 \\
STCG-v1/no-state & 0.798 & 0.646 & 38.6 \\
STCG-v1/state-only & \textbf{0.812} & \textbf{0.673} & \textbf{35.5} \\
STCG-v1/state+contrast & 0.808 & 0.665 & 36.4 \\
\bottomrule
\end{tabular}
\Description{Ablation table comparing DynVAR-Lasso, STCG-v1 without state aggregation, STCG-v1 with state-only aggregation, and STCG-v1 with the optional state-contrast term.}
\end{table}

The ablation in Table~\ref{tab:state-ablation} isolates the mechanism.
Removing state clustering collapses to the dynamic window baseline, while inferred state aggregation improves Edge AUC and SHD.
The contrast term is therefore treated as a negative ablation rather than a core contribution.

\section{Analysis}

\begin{figure}[t]
\centering
\includegraphics[width=\linewidth]{results/figures/stcg_v1_state_diagnostics_seed0.png}
\caption{Window-level graph-state diagnostics for a switching synthetic run. STCG-v1 infers states without regime labels; after label alignment, the only seed-0 mismatch is the window crossing the true switch boundary.}
\Description{Timeline and confusion-matrix diagnostic comparing inferred STCG graph states with hidden switching regimes.}
\label{fig:state-diagnostics}
\end{figure}

The state diagnostic in Figure~\ref{fig:state-diagnostics} checks whether the inferred clusters correspond to meaningful regime changes.
Across five seeds, STCG-v1 infers two states in every run, reaches aligned state accuracy $0.900 \pm 0.091$, and achieves Adjusted Rand Index $0.595 \pm 0.370$.
Lower ARI cases are mainly caused by windows that straddle the true switching boundary.

Robustness sweeps over observation noise, sequence length, and edge density further clarify the operating regime.
STCG-v1 obtains the best Edge AUC in 10 of 12 sweep settings over twenty seeds.
It is best at all tested observation-noise levels and at moderate to dense graph settings, but short sequences of 300 and 600 timesteps still favor DynVAR-Lasso.
The main failure mode is therefore limited temporal evidence: if a window has too few samples or mixes regimes heavily, state inference becomes less reliable.
The switch-difficulty sweep makes this boundary explicit.
When the regime switch period is 150 timesteps with 300-step windows, oracle window purity falls to $0.500$ and STCG-v1 ties DynVAR-Lasso on Edge AUC ($0.733$).
At a 600-step switch period, oracle purity rises to $0.929$, STCG-v1 reaches Edge AUC $0.847$, and the oracle upper bound is $0.849$ versus DynVAR-Lasso $0.835$.

The real-data results are intentionally framed as probes rather than a leaderboard.
They ask whether graph-derived structure transfers to forecasting, operating-state diagnosis, or future transition warning under chronological evaluation.
This framing prevents a forecasting gain from being treated as the only success criterion and, equally, prevents negative forecasting results from being hidden when graph states still carry diagnostic information.

\subsection{Limitations}

The current evidence has four important boundaries.
First, global VAR-Lasso, Granger-F, and PCMCI-ParCorr remain stronger on stationary and nonlinear VAR, so STCG-v1 should be presented as a dynamic-structure method rather than a universal temporal discovery method.
Second, pair-level AUC and lag accuracy improve in mean on switching VAR but are not consistently statistically significant under the current twenty-seed paired $t$ and Wilcoxon comparisons.
Third, real and smoke-test forecasting gains are limited: hard masks do not beat dense or persistence baselines, the tuned AReM STCG-v1 gain over DynVAR-Lasso is promising but not statistically confirmed, UCI HAR does not reproduce that gain, and HydraulicSystems validates operating-condition state alignment without a forecasting gain.
The state-diagnosis harness includes both label-stratified and forward-only variants, so it supports graph-state interpretability and condition-diagnosis evidence rather than a full deployment claim.
The forward-only cooler result reduces one leakage concern but has a single-condition future test segment, and the multi-class valve forward check is weak for compact STCG states.
The transition-warning pilot is future-facing and includes validation-selected STCG warning ablations, but STCG-v1 warning features are mixed and do not yet replace the simpler Ridge graph or sensor baselines on the strongest AP settings.
Fourth, the current implementation is sparse reconstruction plus graph-state aggregation; although the synthetic recovery loop now includes PCMCI-ParCorr, broader neural self-supervised claims and kernel/nonlinear PCMCI variants require additional implemented objectives and baselines.

\section{Conclusion}

This paper establishes a reproducible synthetic recovery loop for self-supervised lag-aware dynamic graph discovery.
The strongest current result is STCG-v1 on switching-regime systems, where latent graph-state aggregation improves edge recovery and produces interpretable states aligned with hidden regimes.
The real-data story is now sharper: Hydraulic Systems shows that graph-derived states and graph-score features can support operating-condition diagnosis and some future transition warnings even when one-step RMSE does not improve, while AReM and UCI HAR keep the transfer claim bounded.
Overall, the evidence supports STCG-v1 as a dynamic-structure discovery method for regime-changing temporal systems, with real-data warning and diagnosis experiments serving as bounded diagnostic support rather than forecasting state-of-the-art claims.

\paragraph*{Artifact Availability.}
\begin{sloppypar}
An anonymous reviewer artifact with code, configurations, result tables, figures, and reproduction instructions is available at:
\href{https://github.com/artifact-lab/stcg-anonymous-artifact}{\url{https://github.com/artifact-lab/}}\\
\href{https://github.com/artifact-lab/stcg-anonymous-artifact}{\url{stcg-anonymous-artifact}}.
\end{sloppypar}

\section*{GenAI Usage Disclosure}

Generative AI tools were used to assist with research planning, drafting, code scaffolding, and language editing. The authors are responsible for all submitted content, experiments, claims, and verification.

\bibliographystyle{ACM-Reference-Format}
\bibliography{paper/references}

\end{document}
